% ==============================================================================
\chapter{Shape Parameterization: Reference Configuration Mapping}
\label{ch:pullback}
% ==============================================================================

\begin{tcolorbox}[colback=lightbg, colframe=primaryblue, title=\textbf{Chapter at a Glance: Reference Diffeomorphic Mapping}, fonttitle=\bfseries]
\begin{description}[leftmargin=!, labelwidth=3.4cm]
    \item[\textbf{Core Objective}] Map parameterized physical domains $\Omega(\boldsymbol{\mu})$ onto a static reference domain $\hat{\Omega}$ to preserve snapshot vector space compatibility.
    \item[\textbf{Mathematical Tool}] Smooth bijective mapping (diffeomorphism) $\boldsymbol{\Psi}(\hat{\vect{x}}; \boldsymbol{\mu}): \hat{\Omega} \to \Omega(\boldsymbol{\mu})$ with Jacobian $J_{\hat{\vect{x}}, \boldsymbol{\mu}} > 0$.
    \item[\textbf{Key Result}] Shifts geometric parameter dependency entirely into integration Jacobians, enabling static basis functions $\hat{\mat{R}}$ and sub-millisecond hyper-reduction.
\end{description}
\end{tcolorbox}

% ------------------------------------------------------------------------------
\section{Coordinate Domain Mapping Architecture}
% ------------------------------------------------------------------------------

To resolve vector space incompatibility across varying geometric parameters $\boldsymbol{\mu}$, we establish a three-domain mapping framework. As illustrated in Figure \ref{fig:pullback_mapping}, the physical coordinates $\vect{x} \in \Omega(\boldsymbol{\mu})$ are expressed as a composition of reference coordinates $\hat{\vect{x}} \in \hat{\Omega}$ and a parameter-dependent displacement field $\vect{d}(\boldsymbol{\mu})$:

\begin{equation}
    \vect{x} = \hat{\vect{x}} + \vect{d}(\boldsymbol{\mu}) = \hat{\boldsymbol{\Psi}}(\boldsymbol{\xi}) + \vect{d}(\boldsymbol{\mu})
\end{equation}

\begin{figure}[htbp]
    \centering
    \makebox[\textwidth][c]{%
        \begin{tikzpicture}[scale=0.9, >=stealth]
            % --- PARAMETRIC DOMAIN (left) ---
            \begin{scope}[xshift=0cm, yshift=0cm]
                \draw[fill=lightbg, draw=primaryblue, thick] (0,0) rectangle (2.5,2.0);
                \foreach \x in {0.5, 1.0, 1.5, 2.0} {
                    \draw[primaryblue!30, thin] (\x, 0) -- (\x, 2.0);
                }
                \foreach \y in {0.5, 1.0, 1.5} {
                    \draw[primaryblue!30, thin] (0, \y) -- (2.5, \y);
                }
                \node[below=2pt] at (1.25, 0) {$\xi$};
                \node[left=4pt] at (0, 1.0) {$\eta$};
                \node[fill=lightbg, inner sep=2pt, rounded corners=1pt] at (1.25, 1.0) {\large $\hat{\Omega}_{\text{param}}$};
                \node[below=14pt, align=center] at (1.25, 0) {\small Parametric Space \\ $\boldsymbol{\xi} = (\xi, \eta)^T$};
            \end{scope}

            % --- REFERENCE CONFIGURATION (top right) ---
            \begin{scope}[xshift=7.0cm, yshift=2.2cm, scale=0.6]
                \draw[fill=gray!5, draw=black, thick]
                (0, -1.0) -- (2.0, -1.0)
                .. controls (2.5, -1.0) and (2.7, -0.4) .. (3.0, -0.4)
                .. controls (3.3, -0.4) and (3.5, -1.0) .. (4.0, -1.0)
                -- (6.0, -1.0) -- (6.0, 1.0) -- (4.0, 1.0)
                .. controls (3.5, 1.0) and (3.3, 0.4) .. (3.0, 0.4)
                .. controls (2.7, 0.4) and (2.5, 1.0) .. (2.0, 1.0)
                -- (0, 1.0) -- cycle;
                
                \foreach \x in {1, 2, 3, 4, 5} {
                    \draw[black!20, thin] (\x, -1.0) -- (\x, 1.0);
                }
                \draw[black!20, thin] (0, -0.5) -- (6.0, -0.5);
                \draw[black!20, thin] (0, 0) -- (6.0, 0);
                \draw[black!20, thin] (0, 0.5) -- (6.0, 0.5);
                
                \fill[black] (4.5, 0.3) circle (2.0pt) node[above] {$\hat{\vect{x}}$};
                \node[fill=gray!5, inner sep=2pt] at (3.0, 0) {\large $\hat{\Omega}$};
                \node[right=15pt, align=left] at (6.0, 0) {\small Reference Domain \\ (Static, $\hat{\vect{x}}$)};
            \end{scope}

            % --- PARAMETERIZED PHYSICAL DOMAIN (bottom right) ---
            \begin{scope}[xshift=7.0cm, yshift=-1.5cm, scale=0.6]
                \draw[fill=lightbg, draw=primaryblue, thick]
                (0, -1.0) -- (2.0, -1.0)
                .. controls (2.4, -1.0) and (2.6, -0.2) .. (3.0, -0.2)
                .. controls (3.4, -0.2) and (3.6, -1.0) .. (4.0, -1.0)
                -- (6.0, -1.0) -- (6.0, 1.0) -- (4.0, 1.0)
                .. controls (3.6, 1.0) and (3.4, 0.2) .. (3.0, 0.2)
                .. controls (2.6, 0.2) and (2.4, 1.0) .. (2.0, 1.0)
                -- (0, 1.0) -- cycle;
                
                \foreach \x in {1, 2, 3, 4, 5} {
                    \draw[primaryblue!20, thin] (\x, -1.0) -- (\x, 1.0);
                }
                \draw[primaryblue!20, thin] (0, -0.5) .. controls (2, -0.5) and (2.5, -0.15) .. (3.0, -0.15) .. controls (3.5, -0.15) and (4, -0.5) .. (6, -0.5);
                \draw[primaryblue!20, thin] (0, 0) -- (6.0, 0);
                \draw[primaryblue!20, thin] (0, 0.5) .. controls (2, 0.5) and (2.5, 0.15) .. (3.0, 0.15) .. controls (3.5, 0.15) and (4, 0.5) .. (6, 0.5);
                
                \fill[primaryblue] (4.5, 0.3) circle (2.0pt) node[above] {$\vect{x}$};
                \node[fill=lightbg, inner sep=2pt] at (3.0, 0) {\large $\Omega(\boldsymbol{\mu})$};
                \node[right=15pt, align=left] at (6.0, 0) {\small Physical Domain \\ (Parameterized, $\vect{x}$)};
            \end{scope}

            % --- MAPPING ARROWS ---
            \draw[->, ultra thick, black!60] (2.7, 1.5) to[out=25, in=195] node[midway, above left, black] {$\hat{\boldsymbol{\Psi}}(\boldsymbol{\xi})$} (6.8, 2.5);
            \draw[->, ultra thick, black!60] (2.7, 0.5) to[out=-25, in=165] node[midway, below left, black] {$\boldsymbol{\Psi}(\boldsymbol{\xi}; \boldsymbol{\mu})$} (6.8, -0.6);
            \draw[->, ultra thick, red!70!black] (8.8, 1.5) to[out=-90, in=90] node[midway, right, red!70!black] {$\vect{d}(\boldsymbol{\mu})$} (8.8, -0.2);
            
            \node[below, font=\small] at (8.8, -2.5) {$\vect{x} = \hat{\vect{x}} + \vect{d} = \hat{\boldsymbol{\Psi}}(\boldsymbol{\xi}) + \vect{d}(\boldsymbol{\mu})$};
        \end{tikzpicture}%
    }
    \caption{Diffeomorphic coordinate mappings between Parametric Space $\hat{\Omega}_{\text{param}}$, static Reference Domain $\hat{\Omega}$, and Parameterized Physical Domain $\Omega(\boldsymbol{\mu})$.}
    \label{fig:pullback_mapping}
\end{figure}

% ------------------------------------------------------------------------------
\section{Continuum Pullback Formulation}
% ------------------------------------------------------------------------------

Let $\boldsymbol{\Psi}(\hat{\vect{x}}; \boldsymbol{\mu}) : \hat{\Omega} \to \Omega(\boldsymbol{\mu})$ be a smooth, orientation-preserving diffeomorphism with positive Jacobian determinant $J_{\hat{\vect{x}}, \boldsymbol{\mu}} = \det \mathbf{J}_{\hat{\vect{x}}} > 0$. Physical field gradients and volume elements pull back to $\hat{\Omega}$ via:

\begin{equation}
    \nabla_{\vect{x}} u = \mathbf{J}_{\hat{\vect{x}}}^{-T} \nabla_{\hat{\vect{x}}} \hat{u}, \quad \mathrm{d}\Omega = J_{\hat{\vect{x}}, \boldsymbol{\mu}} \, \mathrm{d}\hat{\Omega}
\end{equation}

Substituting these transformation laws into the elastodynamic weak form yields the pulled-back variational statement on the static reference configuration $\hat{\Omega}$:

\begin{equation}
    \hat{m}(\ddot{\vect{u}}, \delta\vect{u}; \boldsymbol{\mu}) + \hat{k}(\vect{u}, \delta\vect{u}; \boldsymbol{\mu}) = \hat{f}(\delta\vect{u}; \boldsymbol{\mu}) \quad \forall \delta\vect{u} \in \mathcal{V}_0(\hat{\Omega})
\end{equation}

where the mapped bilinear and linear forms are evaluated as:
\begin{align}
    \hat{m}(\ddot{\vect{u}}, \delta\vect{u}; \boldsymbol{\mu}) &= \int_{\hat{\Omega}} \rho \ddot{\vect{u}} \cdot \delta\vect{u} \, J_{\hat{\vect{x}}, \boldsymbol{\mu}} \, \mathrm{d}\hat{\Omega} \\
    \hat{k}(\vect{u}, \delta\vect{u}; \boldsymbol{\mu}) &= \int_{\hat{\Omega}} \boldsymbol{\varepsilon}_{\hat{\vect{x}}}(\delta\vect{u})^T \mathbf{J}_{\hat{\vect{x}}}^{-1} \mat{D} \mathbf{J}_{\hat{\vect{x}}}^{-T} \boldsymbol{\varepsilon}_{\hat{\vect{x}}}(\vect{u}) \, J_{\hat{\vect{x}}, \boldsymbol{\mu}} \, \mathrm{d}\hat{\Omega} \\
    \hat{f}(\delta\vect{u}; \boldsymbol{\mu}) &= \int_{\hat{\Omega}} \vect{f}_b \cdot \delta\vect{u} \, J_{\hat{\vect{x}}, \boldsymbol{\mu}} \, \mathrm{d}\hat{\Omega} + \int_{\hat{\Gamma}_N} \vect{h}_N \cdot \delta\vect{u} \, J_{\hat{\vect{x}}, \boldsymbol{\mu}, \Gamma} \, \mathrm{d}\hat{\Gamma}
\end{align}

% ------------------------------------------------------------------------------
\section{Naive vs. Pullback Algebraic Paradigms}
% ------------------------------------------------------------------------------

Table \ref{tab:naive_vs_pullback_algebraic} compares the discretization properties resulting from the Naive and Pullback approaches.

\begin{table}[htbp]
    \centering
    \caption{Algebraic and computational comparison between Naive and Pullback discretization.}
    \label{tab:naive_vs_pullback_algebraic}
    \resizebox{\textwidth}{!}{%
    \begin{tabular}{p{3.5cm} p{5.5cm} p{5.5cm}}
        \toprule
        \textbf{Feature} & \textbf{Naive Physical Discretization} & \textbf{Pullback Reference Discretization} \\
        \midrule
        \textbf{Basis Functions} & Parameter-dependent $\mat{R}(\boldsymbol{\mu})$ & Static reference basis $\hat{\mat{R}}$ \\
        \textbf{Vector Space} & Variable DOF count across $\boldsymbol{\mu}$ & Constant reference DOF vector space \\
        \textbf{POD / SVD Basis} & Impossible (snapshot size mismatch) & Native snapshot matrix POD projection \\
        \textbf{Online Assembly} & Full physical remeshing \& re-assembly & Fast Jacobian updates at integration points \\
        \bottomrule
    \end{tabular}%
    }
\end{table}

Under the pullback paradigm, because reference basis functions $\hat{\mat{R}}$ do not depend on parameter $\boldsymbol{\mu}$, all snapshot vectors belong to a single vector space $\mathbb{R}^{N_{dof}}$, making Proper Orthogonal Decomposition (POD) mathematically consistent and exact.

% ------------------------------------------------------------------------------
\section{Chapter Summary}
% ------------------------------------------------------------------------------

\begin{tcolorbox}[colback=lightbg, colframe=accentcyan, title=\textbf{Key Takeaway}, fonttitle=\bfseries]
The diffeomorphic pullback mapping shifts parameter dependencies into the integration Jacobians $\mathbf{J}_{\hat{\vect{x}}}(\boldsymbol{\mu})$ while preserving a constant reference discretization $\hat{\mat{R}}$. This guarantees vector space compatibility across parameter instances, laying the foundation for snapshot POD collection in Chapter \ref{ch:simulation} and ROM hyper-reduction in Chapter \ref{ch:rom}.
\end{tcolorbox}
